\begin{textsample}{POS Dim 3 – 2009\_c6 – Score -3.00 – 2009\_c6/enade\_2009\_com\_soc\_rad\_1bee88f743.txt}  \label{ex:f3_pos_245}
Capítulo 3 Análise Técnica da Prova Este capítulo tem por objetivo apresentar o desempenho dos estudantes de COMUNICAÇÃO SOCIAL – RADIALISMO no ENADE/2009.
Para isso, foram calculadas as estatísticas básicas da prova como um todo, bem como as estatísticas das questões relacionadas à Formação Geral e ao Componente Específico.
Nas Tabelas, são apresentadas as seguintes estatísticas : total da população de presentes, média, erro-padrão da média, desvio-padrão, mediana, nota máxima, nota mínima e coeficiente de assimetria ( ver definição em 1.4.3 ).
Nas estatísticas das questões discursivas, foram detalhadas também as notas nulas : em branco, por protesto ou por insuficiência.
As estatísticas apresentadas neste capítulo contemplam, separadamente, os ingressantes, os concluintes e o total de estudantes, e foram calculadas tendo -se em vista as seguintes agregações : ( a ) Região Geográfica e País como um todo e ( b ) a Categoria Administrativa e a Organização Acadêmica.
As notas brutas das questões discursivas em cada uma das partes, Formação Geral/Componente Específico, foram calculadas a partir da média simples das notas de cada uma das questões correspondentes.
A primeira questão discursiva da Formação Geral foi desconsiderada e, portanto, a nota bruta referente às questões discursivas de Formação Geral foi igual à nota da questão mantida, a segunda.
No Componente Específico, a nota bruta correspondeu à média das notas nas três questões.
As notas brutas das questões discursivas em cada uma das partes, Formação Geral/Componente Específico, foram calculadas a partir dos acertos dos estudantes considerando -se uma regra de três para a conversão para a escala de 0 a 100 descontando -se aquelas questões que foram desconsideradas pela comissão ou que não passaram pelo teste do índice de discriminação ponto-bisserial ( ver seção 1.4.6 para a descrição do índice ).
Questões com índices de discriminação inferior a 0,2 foram excluídas do cômputo da nota do estudante.
As notas brutas de Formação Geral e no Componente Específico para cada um dos estudantes foram calculadas como médias ponderadas das suas notas nas questões objetivas e discursivas, ponderação diferente para cada parte : C\_FG = 0,6 × C\_FG^OBJ + 0,4 × C\_FG^DIS C\_CE = 0,85 × C\_CE^OBJ + 0,15 × C\_CE^DIS, onde C\_FG é a nota bruta da Formação Geral, C\_FG^OBJ é a nota bruta nas questões objetivas da Formação Geral, C\_FG^DIS é a nota bruta nas questões discursivas da Formação Geral, C\_CE é a nota bruta do Componente Específico, C\_CE^OBJ é a nota bruta nas questões objetivas do Componente Específica, e C\_CE^DIS é a nota bruta nas questões discursivas do Componente Específico.
A nota bruta da prova, C\_GER, por sua vez, é também uma média ponderada da nota de cada uma das partes, Formação Geral/Componente Específico : C\_GER = C\_FG × 0,25 + C\_CE × 0,75. 3.1 Estatísticas Básicas da Prova 3.1.1 Estatísticas Básicas Gerais A Tabela 3.1 apresenta as estatísticas básicas da prova por grupo de estudantes.
O total de estudantes inscritos para a prova foi de 2982.
Destes, 79,2 \% compareceram, sendo que a presença foi maior entre concluintes ( 86,7 \% ) do que entre ingressantes ( 71,3 \% ).
A média geral da prova foi de 40,19, sendo que os ingressantes obtiveram média mais baixa, de 38,79, que os concluintes, cuja média foi 41,29.
O desvio-padrão geral foi de 16,4, sendo o do grupo de ingressantes menor ( 15,5 ), indicando que, em contraposição ao grupo dos concluintes ( 17,0 ), aqueles possuíram uma distribuição mais homogênea nas notas.
A nota máxima foi 87,8, obtida por pelo menos um concluinte ao passo que a maior nota obtida por um ingressante foi 84,7.
Os coeficientes de assimetria para os dois grupos considerados, o de concluintes e o de ingressantes, bem como para a população como um todo foram todos negativos, indicando uma maior concentração à esquerda da distribuição e um maior espalhamento à direita.
Uma comparação entre o desempenho de ingressantes e concluintes pode ser realizada a partir dos dados apresentados no Gráfico 3.1.
Cumpre destacar que as notas dos concluintes, de um modo geral, foram mais elevadas do que as dos ingressantes, cuja distribuição de notas está à esquerda dos primeiros.
A ideia de melhor desempenho é reforçada pelas notas medianas dos dois grupos : 43,2 entre os concluintes e 39,7 entre os ingressantes.
Nas faixas de notas até 40 pontos o percentual de ingressantes é maior que o dos concluintes, situação esta que se inverte para as faixas de notas superiores a 40 pontos.
A exceção é o primeiro intervalo de 0 a 10 que inclui uma quantidade maior de concluintes do que de ingressantes.
Constatou -se que aproximadamente 22,6 \% dos ingressantes e 31,9 \% dos concluintes situam -se nas faixas de notas superiores a 50,0 pontos.
O Gráfico 3.2 apresenta informações referentes ao desempenho geral de ingressantes e concluintes, comparando os resultados considerando a Região, Categoria Administrativa e Organização Acadêmica.
Levando -se em conta as notas médias dos estudantes em cada Região, observa -se que, há nos dois grupos, uma variação expressiva entre as Regiões.
As notas médias mais elevadas entre os concluintes ocorreram nas Regiões Sul ( 46,81 ) e Nordeste ( 45,50 ), ambas acima da média nacional ( 41,28 ).
Em relação aos ingressantes, as médias mais altas foram obtidas nas Regiões Centro-Oeste ( 48,74 ), Nordeste ( 42,12 ) e Sul ( 42,35 ) nas demais Regiões a média foi inferior à média nacional ( 38,79 ).
Na Região Norte, foi obtida a média de notas mais baixa entre os ingressantes ( 29,15 ).
Entre os concluintes, a nota mais baixa foi verificada na Região Centro-Oeste ( 38,77 ).
O Centro-Oeste apresentou a média dos ingressantes superior a dos concluintes, fato devido ao maior número de provas com nota zero entre os concluintes.
Este curso é oferecido em instituições federais, estaduais, municipais e privadas.
Levando -se em conta os agrupamentos dos estudantes em categorias administrativas, observa -se que a pontuação média mais elevada entre os concluintes e ingressantes foi encontrada em instituições de origem federal ( respectivamente 42,92 e 44,27 ) – ambas acima da média nacional.
Tendo como foco as notas médias dos estudantes agrupados por Organização Acadêmica, notamos que Universidades apresentam as maiores médias entre os concluintes ( 41,63 ) e Faculdades apresentam as maiores médias entre os ingressantes ( 39,53 ).
Ao passo que as piores notas médias foram observadas nos Centros Universitários para concluintes ( 40,02 ) e nas Universidades para ingressantes ( 38,43 ). 3.1.2 Estatísticas Básicas de Formação Geral A Tabela 3.2 apresenta as estatísticas básicas de Formação Geral por grupo de estudantes.
O total de estudantes inscritos para a prova foi de 2982.
Destes, 79,2 \% compareceram, sendo que a presença foi maior entre concluintes ( 86,7 \% ) do que entre ingressantes ( 71,3 \% ).
A média geral da parte da prova de Componente de Formação Geral foi de 38,82, sendo que os ingressantes obtiveram média mais baixa, 38,32, que os concluintes, cuja média foi 39,22.
O desvio-padrão geral foi de 21,1, do que o do grupo de ingressantes menor ( 20,2 ), indicando que, em contraposição ao grupo dos concluintes ( 21,8 ), aqueles possuíram uma distribuição mais homogênea nas notas.
A nota máxima foi 91,4, obtida por estudantes de ambos os grupos.
Para esta parte, assim como para a nota geral, os coeficientes de assimetria são positivos para todos os grupos.
O Gráfico 3.3 permite a comparação entre concluintes e ingressantes relativamente ao desempenho na parte da prova que aborda a Formação Geral.
Nesta parte, se excluídas as notas nulas, tal como na prova como um todo, também os concluintes apresentaram desempenho superior, embora seja perceptível um maior equilíbrio.
A mediana assim como a média dos concluintes foi maior que a dos ingressantes.
Dos ingressantes, aproximadamente 28,7 \% alcançaram nota superior a 50 ; dentre os concluintes, o percentual dos estudantes nesta faixa de notas é muito semelhante e igual a 32,4 \%.
No Gráfico 3.4 são apresentadas as informações referentes ao desempenho dos ingressantes e dos concluintes na parte de Formação Geral, em diferentes agregações como Região, Categoria Administrativa e Organização Acadêmica.
Considerando -se as notas médias dos estudantes segundo cada Região, observou -se que, tanto no grupo dos ingressantes quanto no grupo dos concluintes, houve diferenças marcantes.
As médias variam, entre os concluintes, de 33,99 na Região Centro-Oeste a 45,83 na Região Sul.
Entre os ingressantes, os valores foram de 37,15 na Região Sudeste a 50,83 na Região Centro-Oeste.
A maior dispersão de notas dos estudantes de uma dada Região foi encontrada no Norte tanto para os concluintes quanto para os ingressantes.
Nas Regiões Norte, Nordeste, Sudeste e Sul, os concluintes apresentaram médias maiores do que os ingressantes, a situação inversa acontecendo no Centro-Oeste.
Todas as Regiões apresentaram concluintes.
Tendo como foco as categorias administrativas, destacou -se que todas as instituições estavam representadas.
As notas médias nas federais foram superiores para ingressantes ( 46,20 ).
As notas médias das demais categorias administrativas foram superiores para concluintes.
Considerando -se o tipo de Organização Acadêmica, somente as Universidades apresentaram o maior valor para concluintes ( 40,02 ).
E as Faculdades apresentaram o maior valor para ingressantes ( 39,11 ), sendo este inclusive superior à nota dos seus concluintes.
Os Centros Federais de Educação Tecnológica e os Institutos Federais de Educação, Ciência e Tecnologia não estavam representados nesta área. 3.1.3 Estatísticas Básicas do Componente Específico A Tabela 3.3 apresenta as estatísticas básicas referentes ao Componente Específico da prova, tendo sua média geral, de 40,69, um pouco superior ao desempenho da prova de COMUNICAÇÃO SOCIAL – RADIALISMO como um todo, cuja média foi 40,19, como apresentado na Tabela 3.1, que por sua vez foi superior ao desempenho na Formação Geral, que teve média igual a 38,82, conforme mostrado na Tabela 3.2.
Os concluintes apresentaram um desempenho médio ( 42,02 ) superior ao dos ingressantes ( 39,00 ), havendo ainda indícios de maior homogeneidade entre as notas médias dos ingressantes do que entre as dos concluintes, como sugeriu o desvio padrão, menor no grupo dos ingressantes do que no dos concluintes.
As notas máximas obtidas por concluintes ( 86,7 ) e ingressantes ( 83,8 ) foram diferentes.
As distribuições de notas para cada grupo se apresentaram com uma inclinação para notas maiores, como evidenciado pelos coeficientes de assimetria negativos.
A situação de assimetria é confirmada pelo Gráfico 3.5 que proporciona uma comparação do desempenho entre concluintes e ingressantes em relação ao Componente Específico da prova.
Claramente a distribuição das notas médias dos ingressantes, à esquerda daquela dos concluintes, evidencia o percentual mais elevados de notas médias nas faixas mais baixas para este grupo.
As medianas das distribuições das notas dos ingressantes e dos concluintes foram 40,0 e 44,2, respectivamente, enfatizando o melhor desempenho dos concluintes nesta parte da prova.
Assim como os Gráficos 3.2 e 3.4, o Gráfico 3.6 apresenta uma comparação dos resultados em relação às Regiões, Categoria Administrativa e Organização Acadêmica, agora levando em conta o desempenho dos ingressantes e dos concluintes no Componente Específico da prova.
Tendo como foco as notas médias dos estudantes, em determinadas Regiões, observa -se que, entre os concluintes, as mais elevadas foram encontradas nas Regiões Sul ( 47,18 ) e Nordeste ( 45,78 ), valores acima da média nacional para esta parte ( 42,02 ).
Entre os ingressantes, as maiores notas ocorreram nas Regiões Centro-Oeste ( 48,10 ), Sul ( 42,59 ) e Nordeste ( 41,49 ), todas acima da média nacional ( 38,99 ).
A Região Sudeste apresenta os contingentes mais homogêneos com respeito à componente específica, tanto entre os ingressantes ( 0,51 ) como entre os concluintes ( 0,55 ).
Novamente, para a análise segundo a Categoria Administrativa, todas as entidades foram representadas.
Observa -se comportamento diverso àquele da Formação Geral, pois entre os concluintes as médias mais elevadas na Formação Geral foram entre os de origem municipal e no Componente Específico passam a estar entre os de origem federal ( 43,13 ).
Já para os ingressantes, o comportamento foi semelhante, ao da parte da Formação Geral e da parte do Componente Específico, mantendo as médias mais elevadas entre os de origem federal ( 43,67 ).
Quanto à Organização Acadêmica, o comentário é o mesmo já feito para a nota bruta e referente à Formação Geral : somente as Universidades, Centros Universitários e Faculdades foram representados neste curso.
Considerando -se o tipo de Organização Acadêmica, entre os concluintes, Universidades apresentaram os maiores valores, seguidas das Faculdades e dos Centros Universitários.
E entre os ingressantes, Faculdades apresentaram os maiores valores, seguidas dos Centros Universitários e das Universidades. 3.2 Análise das Questões Objetivas 3.2.1 Formação Geral A Tabela 3.4 apresenta as estatísticas básicas relativas às oito questões objetivas da Formação Geral dos estudantes.
Adiantamos que a primeira questão foi desconsiderada pelo critério da correlação ponto-bisserial.
Como pode ser observado na Tabela 3.4, os estudantes concluintes obtiveram desempenho superior, com nota média de 49,49, em relação aos estudantes ingressantes cuja nota média foi 48,73.
A variabilidade entre os concluintes foi um pouco maior do que entre os ingressantes, como indicaram os desvios-padrão de 22,0, para o grupo de ingressantes e 24,2 para o grupo dos concluintes.
A amplitude total das notas foi a mesma para ambos os grupos : nota mínima de 0,00 e máxima de 85,7.
A mediana foi igual para ambos os grupos : 57,1.
A assimetria, medida pelo coeficiente de assimetria ( skewness ) foi igual para ambos os grupos ( -0,5 ), ainda que ambos fossem negativos indicando uma assimetria para a direita, fato reforçado pelos valores menores das médias dos dois grupos quando comparados à mediana correspondente.
A Tabela 3.5 apresenta a classificação das questões objetivas da Formação Geral, segundo o índice de facilidade ( questões em vermelho foram desconsideradas pelo critério ponto-bisserial ).
Como as questões foram comuns a todas as áreas, a tabela apresenta a classificação considerando -se os concluintes de todas as áreas e aqueles do curso de COMUNICAÇÃO SOCIAL – RADIALISMO.
De acordo com os índices obtidos, as questões objetivas da prova foram assim avaliadas para a área : das oito questões, duas tiveram entre 60 \%, inclusive, e 85 \% de acertos e foram classificadas como fáceis ( somente uma para o total de concluintes ) ; outras três foram consideradas medianas, alcançando entre 40 \%, inclusive, e 60 \% de acertos ( três para o total de concluintes ) ; duas obtiveram entre 15 \%, inclusive, e 40 \% de acertos, sendo consideradas difíceis ( três para o total de concluintes ).
Por fim, apenas uma foi classificada como muito difícil tento obtido percentual menor do que 15 \% de acertos ( uma para o total de concluintes ), coincidentemente a questão desconsiderada pelo critério ponto-bisserial.
Nenhuma questão foi classificada na categoria muito fácil.
A Tabela 3.6 descreve os resultados das análises das questões objetivas relativas à Formação Geral, segundo o poder de discriminação, utilizando -se para tal o índice de discriminação ponto-bisserial.
Como as questões foram comuns a todas as áreas, a tabela apresenta a classificação considerando -se os concluintes de todas as áreas e aqueles do curso de COMUNICAÇÃO SOCIAL – RADIALISMO.
A desconsideração da questão segundo este critério considerou o universo de todos os concluintes.
Cinco das oito questões apresentaram índices iguais ou acima de 0,40 ( nenhuma para a população de concluintes como um todo ), sendo, portanto, classificadas como muito boas nesta característica, para este grupo de estudantes, nenhuma foi classificada como boa com índice entre 0,30, inclusive, e 0,40 ( três para os concluintes como um todo ) e duas foram classificadas como médias, com índice entre 0,20, inclusive, e 0,30 ( quatro para os concluintes como um todo ).
Uma das questões foi classificada como fraca tanto para a população de concluintes como um todo como para a de concluintes de COMUNICAÇÃO SOCIAL – RADIALISMO e foi desconsiderada pelo critério. total de concluintes ) ; outras três foram consideradas medianas, alcançando entre 40 \%, inclusive, e 60 \% de acertos ( três para o total de concluintes ) ; duas obtiveram entre 15 \%, inclusive, e 40 \% de acertos, sendo consideradas difíceis ( três para o total de concluintes ).
Por fim, apenas uma foi classificada como muito difícil tento obtido percentual menor do que 15 \% de acertos ( uma para o total de concluintes ), coincidentemente a questão desconsiderada pelo critério ponto-bisserial.
Nenhuma questão foi classificada na categoria muito fácil.
A Tabela 3.6 descreve os resultados das análises das questões objetivas relativas à Formação Geral, segundo o poder de discriminação, utilizando -se para tal o índice de discriminação ponto-bisserial.
Como as questões foram comuns a todas as áreas, a tabela apresenta a classificação considerando -se os concluintes de todas as áreas e aqueles do curso de COMUNICAÇÃO SOCIAL – RADIALISMO.
A desconsideração da questão segundo este critério considerou o universo de todos os concluintes.
Cinco das oito questões apresentaram índices iguais ou acima de 0,40 ( nenhuma para a população de concluintes como um todo ), sendo, portanto, classificadas como muito boas nesta característica, para este grupo de estudantes, nenhuma foi classificada como boa com índice entre 0,30, inclusive, e 0,40 ( três para os concluintes como um todo ) e duas foram classificadas como médias, com índice entre 0,20, inclusive, e 0,30 ( quatro para os concluintes como um todo ).
Uma das questões foi classificada como fraca tanto para a população de concluintes como um todo como para a de concluintes de COMUNICAÇÃO SOCIAL – RADIALISMO e foi desconsiderada pelo critério.
Na sequência, no Quadro 3.1, estão especificadas, por questão, as habilidades aferidas pelas questões de Formação Geral da prova da COMUNICAÇÃO SOCIAL – RADIALISMO, bem como os resultados da análise relativa aos índices de facilidade e de discriminação de cada questão.
O índice de facilidade para a área de RADIALISMO teve uma variação considerável, oscilando entre 0,08 até 0,78 e o de discriminação ficou entre 0,09 e 0,59.
A questão 8 aferiu as habilidades de “ Ler e interpretar textos ”, “ Analisar e criticar informações ”, “ Extrair conclusões por indução e/ou dedução ”, “ Estabelecer relações, comparações e contrastes em diferentes situações ”, “ Detectar contradições ”, “ Fazer escolhas valorativas avaliando conseqüências ”, “ Questionar a realidade ”, e “ Argumentar coerentemente ” foi a que apresentou, para a população de concluintes de COMUNICAÇÃO SOCIAL – RADIALISMO, maior poder discriminatório, ambas com índice de 0,59.
A questão mais fácil, com 78 \% de acertos, foi também a de número 8.
De outra parte, a questão com menor índice de facilidade e de discriminação em Formação Geral foi a questão 1, que foi desconsiderada, e que deveria aferir as habilidades de “ Analisar informações ” ; “ Extrair conclusões ” e “ Propor soluções para situações-problema ”.
O Gráfico 3.7 mostra, a título de ilustração, o comportamento das respostas à questão 7 de Formação Geral para os estudantes concluintes da área de RADIALISMO como função da nota geral na prova.
Trata -se de uma questão considerada de dificuldade fácil para os concluintes de COMUNICAÇÃO SOCIAL – RADIALISMO e de dificuldade média para o conjunto de concluintes de todas as áreas.
Neste Gráfico, cada uma das cinco curvas, representou o percentual de respostas em determinada alternativa da questão, em função do escore total dos estudantes concluintes da área de RADIALISMO.
A curva azul marinho corresponde à alternativa A, a correta para esta questão.
Assim, observa -se que entre os estudantes com notas mais baixas – desempenho geral mais fraco na prova como um todo –, observa -se, além da escolha da alternativa correta, outras alternativas incorretas, especialmente a alternativa C – curva em vermelho seguida da alternativa B – curva verde claro.
A proporção de estudantes que selecionaram a resposta correta A, aumentou com o desempenho deste grupo de estudantes chegando ao acerto de 100 \% dos estudantes com notas pelo menos 75,0.
Esta análise permitiu verificar como as questões discriminou entre os grupos de desempenho.
Os Gráficos relativos às demais questões de Formação Geral constam do Anexo I. 3.2.2 Componente Específico A Tabela 3.7 apresenta as estatísticas básicas em relação às questões objetivas do Componente Específico.
Como podem ser observados, os estudantes concluintes obtiveram desempenho superior, mas por uma pequena margem com uma média de 45,15 pontos, em relação aos estudantes ingressantes com uma média de 42,36 pontos.
Os desvios-padrão indicaram que a variabilidade entre os concluintes foi maior ( 18,5 ) do que entre os ingressantes ( 17,2 ).
Os coeficientes de assimetria para os dois grupos foram negativos indicando uma inclinação para a direita.
A assimetria foi, porém, maior entre os concluintes.
A análise das questões objetivas do Componente Específico em COMUNICAÇÃO SOCIAL – RADIALISMO, em relação ao grau de facilidade da prova ( ver Tabela 3.8 ), mostrou que, das 27 questões, uma encontraram -se na categoria difícil ( entre 16 \% e 40 \%, não incluso, de acertos ), treze são de nível médio ( entre 41 \% e 60 \%, não incluso, de acertos ), oito estão na categoria fácil ( entre 61 \% e 85 \%, não incluso, de acertos ) e três tiveram a classificação muito fácil ( a partir de 85 \% de acertos ).
Duas questões tiveram percentual de acertos menor do que 15 \% ( classificação muito difícil ) e três foram canceladas pela comissão.
Das 27 questões objetivas relativas ao Componente Específico ( ver Tabela 3.9 ), nove obtiveram o índice de discriminação muito bom, duas obtiveram o índice de discriminação bom e sete obtiveram nível de discriminação médio.
O índice fraco que eliminaria a questão, foi alcançado por sete questões, sendo que duas questões foram canceladas pela comissão.
Na sequência, no Quadro 3.2, estão especificadas, por questão, as habilidades aferidas pelas questões do Componente Específico da prova, bem como os resultados da análise relativa aos índices de facilidade e de discriminação de cada questão.
Das 27 questões objetivas relativas ao Componente Específico ( ver Tabela 3.9 ), nove obtiveram o índice de discriminação muito bom, duas obtiveram o índice de discriminação bom e sete obtiveram nível de discriminação médio.
O índice fraco que eliminaria a questão, foi alcançado por sete questões, sendo que duas questões foram canceladas pela comissão.
Na sequência, no Quadro 3.2, estão especificadas, por questão, as habilidades aferidas pelas questões do Componente Específico da prova, bem como os resultados da análise relativa aos índices de facilidade e de discriminação de cada questão.
O Gráfico 3.8 mostra, a título de ilustração, o comportamento da questão 11 do Componente Específico para os estudantes concluintes da área de RADIALISMO como função da nota geral na prova.
Trata -se de uma questão considerada fácil para os concluintes de COMUNICAÇÃO SOCIAL – RADIALISMO.
O poder de discriminação foi muito bom para os estudantes da área.
Neste Gráfico, cada uma das cinco curvas, representa o percentual de respostas em determinada alternativa da questão, em função da nota dos estudantes concluintes de RADIALISMO.
A curva verde claro correspondeu à alternativa B, a correta para esta questão.
Assim, observou -se que entre os estudantes com escores mais baixos – desempenho geral mais fraco na prova como um todo –, a tendência foi a escolha de uma alternativa incorreta, a alternativa A – curva azul marinho ou a alternativa E – curva roxa.
A proporção de estudantes que selecionaram a resposta correta aumentou com o desempenho deste grupo de estudantes chegando ao acerto de 100 \% dos estudantes com notas pelo menos 75.
Esta análise permitiu verificar como a questão discriminou entre os dois grupos de desempenho.
Os Gráficos relativos às demais questões do Componente Específico constam do Anexo I. 3.3 Análise das Questões Discursivas 3.3.1 Formação Geral A Tabela 3.10 apresenta as estatísticas básicas da prova por grupo de estudantes.
A análise dos resultados de desempenho dos estudantes nas questões discursivas que são apresentados na tabela 3.10 e no gráfico 3.9 mostrou que as notas foram mais baixas no conjunto dessas questões que no das questões objetivas.
O total de estudantes inscritos para a prova foi de 2.982.
Destes, 79,2 \% compareceram, sendo que a presença foi maior entre concluintes ( 86,7 \% ) do que entre ingressantes ( 71,3 \% ).
A média geral da prova foi de 23,43, sendo que os ingressantes obtiveram média mais baixa, de 22,81, que os concluintes, cuja média foi 23,91.
O desvio-padrão geral foi de 32,3, sendo o do grupo de ingressantes menor ( 31,2 ), em relação ao grupo dos concluintes ( 33,1 ), aqueles possuíram uma distribuição mais homogênea nas notas.
A nota máxima foi 100,0, atingida por concluintes e ingressantes.
Os coeficientes de assimetria para os dois grupos considerados, o de concluintes e o de ingressantes, bem como para a população como um todo são todos positivos, indicando uma maior concentração à esquerda da distribuição.
Para todos os grupos mais de metade dos estudantes presentes obtivera nota zero nas questões discursivas da Formação Geral.
Das 1.509 provas com nota zero, 44 ( 2,9 \% ) continham um protesto, 498 ( 33,0 \% ) estavam em branco e 967 ( 64,1 \% ) foram respondidas, porém de forma inadequada.
Entre os concluintes, a maioria ( 62,1 \% ) dentre aqueles com nota zero apresentou a prova respondida de forma inadequada.
Entre os ingressantes 66,6 \% das notas nulas decorreram de correção.
Uma comparação entre o desempenho de ingressantes e concluintes pode ser realizada a partir dos dados apresentados no Gráfico 3.9.
As notas dos concluintes, de um modo geral, foram ligeiramente mais elevadas do que as dos ingressantes, cuja distribuição de notas está à direita dos primeiros, embora a proporção de notas nulas entre os concluintes ( 61,6 \% ) seja maior do que entre os ingressantes ( 59,8 \% ).
A média foi 0,0 tanto para ingressantes como para concluintes.
Na sequência, discorrer -se-á, especificamente, sobre uma das questões discursivas de Formação Geral, a número 10 já que a número 9 foi desconsiderada.
Em especial, serão analisados os conteúdos englobados na questão e os resultados obtidos pelos estudantes.
Ao final, são apresentados os comentários da banca corretora a respeito do observado nas respostas dos estudantes, suas impressões e conclusões.
Cumpre esclarecer que, tendo em vista que as questões discursivas de Formação Geral foram padronizadas, ou seja, constam de todas as provas, os comentários da banca foram os mesmos para todas as áreas acadêmicas, eis que direcionados a todos os estudantes que participaram do ENADE/2009.
A seguir, serão analisados os desempenhos de ingressantes e concluintes da área de COMUNICAÇÃO SOCIAL – RADIALISMO na questão discursiva de Formação Geral, não desconsiderada, do ENADE/2009, comparando os resultados obtidos com a habilidade exigida na questão. 3.3.1.1 Análise da Questão Discursiva 9 de Formação Geral DESCONSIDERADA 3.3.1.2 Análise da Questão Discursiva 10 de Formação Geral A Questão Discursiva nº 10 ousou avaliar as habilidades : • Ler e interpretar textos ; • Analisar e criticar informações ; • Extrair conclusões por indução e/ou dedução ; • Estabelecer relações, comparações e contrastes em diferentes situações ; • Fazer escolhas valorativas avaliando conseqüências ; • Questionar a realidade ; • Argumentar coerentemente.
O Quadro 3.3 apresenta as habilidades/competências, temas predominantes e estatísticas básicas da questão 10 de Formação Geral por grupo de estudantes.
Cumpre notar que como a questão 9 foi desconsiderada, a nota referente às questões discursivas da Formação Geral é idêntica à nota da questão 10.
A média geral desta questão foi de 23,43, sendo que os ingressantes obtiveram média mais baixa, de 22,81, que os concluintes, cuja média foi 23,91.
O desvio-padrão geral foi de 32,3, em relação ao grupo de ingressantes menor ( 31,2 ), indicando que, em contraposição ao grupo dos concluintes ( 33,1 ), aqueles possuíram uma distribuição mais homogênea nas notas.
A nota máxima foi 100,0, atingida por concluintes e ingressantes.
Os coeficientes de assimetria para os dois grupos considerados, o de concluintes e o de ingressante, bem como para a população como um todo foram todos positivos, indicando uma maior concentração à esquerda da distribuição.
Para todos os grupos, mais de metade dos estudantes presentes obtiveram nota zero nesta questão discursiva de Formação Geral.
Das 1.509 provas com nota zero, 44 ( 2,9 \% ) continham um protesto, 498 ( 33,0 \% ) estavam em branco e 967 ( 64,1 \% ) foram respondidas, porém de forma inadequada.
Entre os concluintes, a maioria ( 62,1 \% ) dentre aqueles com nota zero estava com as provas respondidas de forma inadequada.
Entre os ingressantes, 66,6 \% das notas nulas(zero) decorreram do desempenho do próprio estudante verificadas na correção.
Uma comparação entre o desempenho de ingressantes e concluintes para esta questão foi exatamente igual à já realizada para a nota discursiva de Formação Geral a partir dos dados apresentados no Gráfico 3.9. 3.3.1.3 Comentários sobre as respostas à Questão Discursiva 10 Apresentação A questão 10 de Formação Geral apresentou um questionamento capaz de analisar a habilidade do estudante em reconhecer o que é um direito social, optar por um direito social vigente que possa diminuir as desigualdades de renda familiar no Brasil e discutir os argumentos em favor do direito social escolhido como suporte à resposta.
Para análise das respostas apresentadas pelos estudantes, houve uma leitura inicial bastante criteriosa.
Foram considerados para efeito de pontuação, os seguintes argumentos : • Acesso à educação pública, gratuita e de qualidade, garantindo ao cidadão melhores condições de vida ; • Assistência à saúde, em seu contexto mais amplo, garantindo os cuidados essenciais de saúde, baseados em técnicas apropriadas ao alcance de todas as pessoas, contando com a participação da população ; • Permanência do estudante na escola, em todos os níveis escolares ( da educação infantil à educação superior, possibilitando a qualificação profissional ), o que garante melhores condições de trabalho e remuneração.
Considerou -se a qualificação e o desempenho docente, desenvolvendo nos estudantes melhores expectativas de aprendizagem e interesse pelo conteúdo da escola.
Os professores desenvolvendo metodologias adequadas e conhecimento pleno quanto a questões de aprendizagem ; • Condições dignas de trabalho, com remuneração que garanta qualidade de vida.
Foram consideradas as condições de trabalho, a preparação profissional para o trabalho e os cursos de aperfeiçoamento, oportunizando momentos de aprendizagem e desenvolvimento, a partir de situações desencadeadoras ; • Ser proprietário do imóvel em que reside, reduzindo os gastos com aluguel, garantindo melhores condições de vida, equilíbrio financeiro e construção do bem-estar familiar.
Desempenho dos estudantes na questão É importante salientar que havia respostas que desenvolveram adequadamente o questionamento apresentado, com fundamentação baseada em argumentos coerentes.
Contudo, em um número considerável de respostas, havia apenas um argumento e não dois, conforme fora solicitado.
Muitos estudantes informaram sobre problemas da educação, de um modo geral.
Algumas observações mencionadas pelos corretores durante o processo de correção foram : • há preocupação com uma educação de qualidade, com sugestões de cursos profissionalizantes e métodos preparatórios de aulas. • a educação é revelada como de baixa qualidade, exaltando -se treinamentos e cursos de atualização para a classe docente. • a saúde como essência para a vida.
O SUS é colocado como parte integrante de direito social, garantindo às pessoas fácil acesso. • oportunidade de emprego como forma de ascensão social, direito, cidadania. • os políticos brasileiros como causadores do despreparo educacional.
As respostas analisadas apresentaram norma culta, coesão e coerência razoáveis, que garantiram a compreensão da mensagem desenvolvida pelos estudantes.
Os erros ortográficos foram poucos e não comprometeram a discussão sobre o assunto e sobre os pontos de vista abordados.
Quanto à coerência, houve um desvio quanto ao objetivo da questão, já que alguns estudantes, através dos seus argumentos retrataram reivindicações, reclamações e situações de problemas enfrentados pelos brasileiros, citando golpes políticos.
Porém, mesmo quando se tratava desta avaliação, os estudantes foram autênticos através das explicações.
Houve erros quanto à concordância verbal e nominal.
Houve, também, dificuldade em utilizar os conectores corretamente.
Alguns estudantes, por falta de ideias, apresentaram em suas respostas redundância e prolixidade.
O desempenho foi considerado satisfatório, pois a maioria dos estudantes soube responder ao questionamento.
Houve domínio do conteúdo, pois as ideias, os argumentos e a originalidade corresponderam ao esperado. 3.3.2 Componente Específico Nas questões discursivas do Componente Específico – questões de números 38, 39 e 40 –, a diferença de desempenho entre os grupos não foi muito acentuada, inclusive menor do que para as questões discursivas de Formação Geral.
A Tabela 3.11 contém as estatísticas da média destas questões.
Observa -se que no grupo dos concluintes, a nota média alcançada foi 24,61, enquanto que no grupo dos ingressantes foi 20,26, ambas menores do que as obtidas nas questões discursivas de Formação Geral, que foram 38,3 para os ingressantes e 39,2 para os concluintes ( Tabela 3.10 ).
As médias foram também menores dos que as obtidas para as questões objetivas do Componente Específico ( Tabela 3.7 ) : respectivamente 42,4 e 45,2 para ingressantes e concluintes.
Os valores positivos para a medida de assimetria indicaram uma distribuição com cauda para a direita.
A análise de cada uma destas questões será feita nas subseções seguintes.
Dos 2.982 presentes, 43,6 \% obtiveram nota zero.
A tabela 3.12 apresenta um detalhamento da situação destes estudantes.
Um pouco menos de 1/3 ( 28,6 \% ) não fez nenhuma das 3 questões.
Esta proporção foi ligeiramente maior entre os concluintes ( 31,0 \% ) do que entre os ingressantes ( 25,7 \% ).
A combinação de questões em branco e preenchidas correspondeu a 55,3 \% das provas. 1 \% no total das provas correspondeu a 3 questões com protesto.
O Gráfico 3.10 apresenta a distribuição das notas nas questões discursivas do Componente Específico por grupo de estudantes.
A moda se situou no primeiro intervalo, aquele com valores entre 0 e 10 incluindo os extremos do intervalo.
A impressão de assimetria informada pelos coeficientes de assimetria positivos ( Tabela 3.11 ) foi reforçada no gráfico pela concentração no primeiro intervalo. 3.3.2.1 Análise da Questão Discursiva 38 do Componente Específico Questão Desconsiderada. 3.3.2.2 Análise da Questão Discursiva 39 do Componente Específico A questão 39, cujas habilidades aferidas encontram -se descritas no Quadro 3.4, que ainda contém outras informações, seguido da Tabela 3.13 com a distribuição das notas nulas e do Gráfico 3.11, com a distribuição das notas.
Esta questão foi aquela em que os grupos de estudantes observados obtiveram o melhor desempenho entre as questões discursivas do Componente Específico.
A nota média do grupo dos ingressantes foi 27,78, e a dos concluintes 32,11.
O percentual de notas zero obtidas, o mais baixo das questões discursivas do Componente Específico, nos dois grupos, foi de 57,5 \% entre os ingressantes e 54,8 \% entre os concluintes.
Das provas com notas nulas nesta questão, um pouco menos de 1/3 foi deixado em branco : 27,4 \% entre os ingressantes e 33,3 \% entre os concluintes ( ver Tabela 3.14 ).
Entre os com nota nula, uma pequena minoria ( 2,5 \% ) utilizou a questão para protestar.
Apesar do baixo desempenho, em ambos os grupos a nota máxima de 100,0 pontos foi alcançada.
Os coeficientes de assimetria positivos indicaram uma inclinação para a esquerda e cauda para a direita.
O Gráfico 3.11 apresenta a distribuição das notas na questão 39 discursiva do Componente Específico por grupo de estudantes.
A moda se situou no primeiro intervalo, aquele com valores entre 0 e 10 incluindo os extremos do intervalo.
A impressão de assimetria informada pelos coeficientes de assimetria positivos ( Quadro 3.4 ) foi reforçada no gráfico pela concentração no primeiro intervalo.
A distribuição apresentou outros dois máximos locais nos grupos 60 a 70 e 90 a 100, tanto para ingressantes como para concluintes. 3.3.2.4 Comentários sobre as respostas à Questão Discursiva 39 Ao qualificar as ações de marketing como estratégias, há uma simplificação conceitual que neste caso, favoreceu os estudantes no descritivo da resposta demandada pela pergunta.
A perspectiva do consumo incentivado não encontrou barreiras para a maioria dos estudantes inspirar -se e sugerir ações estimuladoras de interatividade comercial.
A maioria dos estudantes ao listar ideias de concursos, prêmios e distinções, projetaram seus anseios e tiveram boas condições de compreensão sobre as expectativas de resultado.
Vive -se em tempos de ampliar a fidelidade comercial e há uma constante troca de mecanismos emocionais de pertencimento simbólico.
O livre trânsito vivencial nestes conceitos facilitou a identificação de respostas objetivas e eficazes para um resultado de médio para bom, nestas questões.
O desempenho dos estudantes foi satisfatório e outras propostas, além das apresentadas na chave de resposta foram consideradas como : as avaliações de conteúdos por pesquisa, o investimento em divulgação do programa e a participação de artistas, entre outras respostas consideradas coerentes pela banca, como instrumento para a ampliação do número de audiência do programa. 3.3.2.5 Análise da Questão Discursiva 40 do Componente Específico A questão 40, cujas habilidades aferidas encontram -se descritas no Quadro 3.5, que ainda contém outras informações, seguido da Tabela 3.14 e do Gráfico 3.12, foi aquela em que os grupos de estudantes obtiveram desempenho melhor entre as questões discursivas do Componente Específico.
A nota média do grupo dos ingressantes foi 22,84, e a dos concluintes 32,14.
O percentual de notas nulas nos grupos foi : 64,7 \% entre os ingressantes e 58,4 \% entre os concluintes.
Aqui, a maioria das provas com nota zero decorre da questão ter sido respondida de forma não satisfatória : 71,3 \% entre os ingressantes e 68,3 \% entre os concluintes ( ver Tabela 3.14 ).
Uma pequena minoria ( 1,9 \% ), dentre aqueles com nota nula, utilizou a questão para protestar.
Em ambos os grupos, a nota máxima de 100,0 pontos foi alcançada.
Os coeficientes de assimetria positivos indicaram uma inclinação para à esquerda.
O Gráfico 3.12 apresenta a distribuição das notas na questão 40 discursiva do Componente Específico por grupo de estudantes.
A moda se situou no primeiro intervalo, aquele com valores entre 0 e 10 incluindo os extremos do intervalo.
A impressão de assimetria informada pelos coeficientes de assimetria positivos ( Quadro 3.5 ) é reforçada no gráfico pela concentração no primeiro intervalo.
A distribuição apresentou outro máximo local no último intervalo, para ingressantes e concluintes. 3.3.2.6 Comentários sobre as respostas à Questão Discursiva 40 Por ser uma questão técnica, em que os estudantes não tiveram conhecimento específico a respeito dos equipamentos necessários para respostas no campo do domínio tecnológico, a digressão eleita pela maioria dos estudantes foi a de apontar para melhores alternativas de equipamento digital, o que já garantiu uma pontuação nesta pergunta.
Ao segmento da questão que exigiu conhecimento técnico do estudante, foi relativamente pequeno o número de estudantes que conseguiu concluir toda a resposta desta questão.
Além do padrão de resposta proposto, a banca corretora considerou correta a proposta que se referiu ao treinamento de pessoal, já que são plausíveis e coerentes aos objetivos propostos no enunciado as argumentações apresentadas pelos estudantes.
De uma forma geral, o desempenho dos estudantes da habilitação de Radialismo adequou -se aos parâmetros de conhecimento propostos nas questões.
O convívio permanente com o produto desenvolvido neste campo de conhecimento trouxe para mais perto do estudante, o universo profissional em questão, através da televisão e do rádio.

% matched lemmas: 
\end{textsample}
